% Nouveau EGA V (cote 25): archive cote 25, site page 143 (PDF page 144).
% typescript. Draft compared: none.
% AI-transcribed by Claude Opus 5.5 (claude-opus-5-5, adaptive thinking, effort high), Message Batches API, batch ega25b (msgbatch_01A7n5YuXubUPNK1zsepi8RG), 2026-11-13; not yet checked by eye.
\scanpage{25/143}
\mspage{\hand{4}}

en supposant X de dimension relative d sur Y en chaque point.  Il y a
lieu de  désigner par \xout{XXX$^{\text{xxp}}$} $V(\phi)^{\text{sup sing}}$ le sous-préschéma fermé de
\xout{Xxxxx} $V(\phi)^{\text{sing}}$ \inserted[typed]{donc de X} défini par l'annulation de cette section (d'un Module
inversible cette fois ci),  dont l'ensemble  sous-jacent est ce qu'il
faut. Il y a lieu de résumer cette construction dans une

\uline{Proposition} 6.

Dans le cas général, on ne peut \xout{dire} dire rien de plus précis sur
$V(\phi)^{\text{sing}}$ et $V(\phi)^{\text{sup sing}}$ .  Nous allons  examiner maintenant un cas
particulier\hand{,} intéressant \xout{\unsure{dans}}pour certaines applications. Nous supposerons
que Y est également lisse sur un préschéma S, avec une dimension relative
constante m pour fixer les idées.  De plus, nous supposons que $V(\phi)^{\text{sing}}$,
que nous notons simplement $V^1$ pour simplifier, défini  par l'annulation
de la section $d^1\phi$  du Module $\underline{P}^1_{X/Y}$ \xout{de rang d} loc libre de \xout{rang} rang d+1,
est  lisse sur S de dimension relative  $(m+d) - (d+1) = m-1$ . (NB  Bien
entendu, les notations $V(\phi)^{\text{sing}}$ et $V(\phi)^{\text{sup sing}}$ sont ambigues en ce sens
qu'elles ne font pas intervenir le préschéma de base \unsure{X|Y} auxquelles elles
se rapportent; dans le cas actuel,  il est sous-entendu que c'est Y,
et on notera d'ailleurs \inserted[typed]{qu'il résulte aussitôt des hypothèses} que tout zéro de $\phi$ \inserted[typed]{singulier} est \uline{non singulier relative-}%
\uline{ment à} S ).  Dans cette situation, nous pouvons écrire  le diagramme
suivant \inserted{de Modules loc. libres sur $V^1$} :

\margin{(D)}
\diagram{\begin{tikzcd}
 & & 0 & & \\
 & & \underline{\Omega}^1_{X/Y}\otimes\underline{O}_{V^1} \arrow[u] & & \\
0 \arrow[r] & \underline{P}^{\hand{\vee}1}_{X/Y}\otimes\underline{O}_{V^1} \arrow[r] \arrow[ur, "\hand{\mu}"] & \underline{\Omega}^1_{X/S}\otimes\underline{O}_{V^1} \arrow[r] \arrow[u] & \underline{\Omega}^1_{V^1/S} \arrow[r] & 0 \\
 & \hand{\underline{O}_{V^1}} \arrow[u] \arrow[r] & \underline{\Omega}^1_{Y/S}\otimes\underline{O}_{V^1} \arrow[u] \arrow[ur, "\hand{\nu}"'] & & \\
\hand{L=\omega^{-2}_{X/Y}\otimes\underline{O}_{V^1}} \arrow[ur] & & 0 \arrow[u] & &
\end{tikzcd}}

\margin{143}
